% Intended LaTeX compiler: pdflatex
\documentclass[a4paper,12pt]{article}
\usepackage[utf8]{inputenc}
\usepackage[T1]{fontenc}
\usepackage{graphicx}
\usepackage{longtable}
\usepackage{wrapfig}
\usepackage{rotating}
\usepackage[normalem]{ulem}
\usepackage{amsmath}
\usepackage{amssymb}
\usepackage{capt-of}
\usepackage{hyperref}
\usepackage{\string~/.emacs.d/common}
\author{mahmood}
\date{\today}
\title{formal system}
\hypersetup{
 pdfauthor={mahmood},
 pdftitle={formal system},
 pdfkeywords={},
 pdfsubject={},
 pdfcreator={Emacs 30.0.50 (Org mode 9.6.6)}, 
 pdflang={English}}
\begin{document}

\maketitle
\begin{note}
this document and others can be found on my website\\
\url{https://mahmoodsheikh36.github.io/}
\end{note}
\begin{definition}
a \textbf{formal system} is an abstract structure used for inferring theorems from \href{20230525165114-axiom.org}{axiom} according to a set of \href{20230524202636-rule_of_inference.org}{rules of inference} which are used for carrying out the inference of theorems from axioms.\\[0pt]
\end{definition}
\begin{my_example}
consider the following formal system \(L_\to\):\\[0pt]
it has the following axioms:\\[0pt]
\begin{itemize}
\item A1. \(\alpha \to (\beta \to \alpha)\)\\[0pt]
\item A2. \((\alpha \to (\beta \to \gamma)) \to ((\alpha \to \beta) \to (\alpha \to \gamma))\)\\[0pt]
\item A3. \((\lnot \beta \to \lnot \alpha) \to ((\lnot \beta \to \alpha) \to \beta)\)\\[0pt]
\end{itemize}
we can substitute any propsitions inplace of \(\beta,\alpha,\gamma\) with the \href{20230401203310-logical_connective.org}{connective}s \(\lnot, \to\) only\\[0pt]
the only rule of inference for this system is \href{20230524210722-modus_ponens.org}{mp}\\[0pt]
\end{my_example}
\begin{definition}
the \href{20230525071323-formal_proof.org}{inference process} in the context of a formal system is as follows:\\[0pt]
let \(\Gamma\) be the set of \href{20230524212241-premise.org}{premise}s (propositions), a proof sequence from \(\Gamma\) is the sequence \(a_1,a_2,\dots,a_n\) such that each statement in the sequence is one of 3 types:\\[0pt]
\begin{enumerate}
\item logical axiom\\[0pt]
\item a statement from \(\Gamma\)\\[0pt]
\item is derived from previous statements in the sequence using one of the rules of inference of the system\\[0pt]
\end{enumerate}
\begin{my_example}
premises: \(\alpha \to \beta, \beta \to \gamma\)\\[0pt]
\href{20230524212250-conclusion.org}{conclusion}: \(\alpha \to \gamma\)\\[0pt]
\begin{proof}
\begin{align*}
  1.\ & \alpha \to \beta & P\\
  2.\ & \beta \to \gamma & P\\
  3.\ & (\alpha \to (\beta \to \gamma)) \to ((\alpha \to \beta) \to (\alpha \to \gamma)) & A2\\
  4.\ & (\beta \to \gamma) \to (\alpha \to (\beta \to \gamma)) & A1\\
  5.\ & \alpha \to (\beta \to \gamma) & 2,4,MP\\
  6.\ & (\alpha \to \beta) \to (\alpha \to \gamma) & 5,3,MP\\
  7.\ & \alpha \to \gamma & 1,6,MP
\end{align*}
\end{proof}
\end{my_example}
\begin{my_example}
proof that \(x \to x\) is a \href{20230525062442-tautology.org}{tautology}\\[0pt]
\begin{proof}
\begin{align*}
  1.\ & x \to ((x \to x) \to x) & A1\\
  2.\ & (x \to ((x \to x) \to x)) \to ((x \to (x \to x)) \to (x \to x)) & A2\\
  3.\ & ((x \to (x \to x)) \to (x \to x)) & 1,2,MP\\
  4.\ & x \to (x \to x) & A1\\
  5.\ & x \to x & 4,3,MP
\end{align*}
\end{proof}
\end{my_example}
\begin{my_example}
an extended formal system \(L_2\)\\[0pt]
axioms:\\[0pt]
A. \(a \lor \lnot a\) such that \(a\) is an arbitrary proposition\\[0pt]
inference rules:  the inference rules based on the logical equivalences and derivations that appear on the \href{20230525082155-logic_formulas.org}{formulas page}\\[0pt]
\begin{subquestion}
prove \(\{D \to B, B \to A, D\} \vdash A\)\\[0pt]
premises: \(D \to B, B \to A, D\)\\[0pt]
conclusion: \(A\)\\[0pt]
\begin{proof}
\begin{align*}
  1.\ & D & P\\
  2.\ & D \to B & P\\
  3.\ & B & 1,2,I11\\
  4.\ & B \to A & P\\
  5.\ & A & 3,4,I11\\
\end{align*}
\end{proof}
\end{subquestion}
\begin{subquestion}
prove using the inference process that\\[0pt]
\begin{itemize}
\item we can infer "joseph doesnt speak truthfully" from the assumptions "joseph doesnt speak truthfully or joseph is lazy" and "joseph isnt lazy"\\[0pt]
\item we can infer \(G \land \lnot E\) from the assumptions \(A \to G, D \land \lnot E, D \to B, B \land A\)\\[0pt]
\item we can infer \(A\) from \(R \to (S \lor A), (P \lor Q) \to R, Q \land (\lnot S)\)\\[0pt]
\end{itemize}
\end{subquestion}
\end{my_example}
\end{definition}
sometimes we might want to prove a conclusion \(C\) with an empty set of premises: \(\vdash C\). an argument of this form is valid if \(C\) is a tautology. meaning \(C\) is a tautology iff we can prove that its always true without the need for assumptions: for example \(\vdash (p \land \lnot p) \to r\):\\[0pt]
\begin{align*}
  1.\ & p \lor \lnot p & A \quad (\text{axiom})\\
  2.\ & \lnot p \lor p & 1,E10\\
  3.\ & \lnot p \lor \lnot (\lnot p) & 2,E9\\
  4.\ & \lnot (p \land \lnot p) & 3,E17\\
  5.\ & \lnot (p \land \lnot p) \lor r & 4,I1\\
  5.\ & (p \land \lnot p) \to r & 5,E20
\end{align*}
\begin{lemma}
if \(\vdash \alpha\) then \(\alpha\) is a tautology, i.e. a proposition that can be formally proved without assumptions is a tautology\\[0pt]
\end{lemma}
\begin{lemma}
if \(\Gamma \vdash \alpha\) then \(\Gamma \implies \alpha\)\\[0pt]
meaning that every \href{20230401175139-statement.org}{proposition} that has a \href{20230525071323-formal_proof.org}{formal proof} in the set of premises \(\Gamma\) logically follows from \(\Gamma\)\\[0pt]
\end{lemma}
\begin{definition}
a system is \textbf{complete} if for every set of assumptions \(\Gamma\), if \(\Gamma \implies \alpha\) then \(\Gamma \vdash \alpha\), i.e. we can \href{20230525071323-formal_proof.org}{formally prove} every true argument with it\\[0pt]
\begin{lemma}
\((\Gamma \implies \alpha) \iff (\Gamma \vdash \alpha)\)\\[0pt]
\end{lemma}
\begin{lemma}
let \(\Gamma\) be a set of assumptions, let \(\alpha,\beta\) be statements, then \(\Gamma \cup \{\alpha\} \vdash \beta\) iff \(\Gamma \vdash (\alpha \to \beta)\)\\[0pt]
\end{lemma}
\begin{lemma}
let \(\Gamma\) be a set of assumptions, let \(\alpha,\beta\) be statements, then \(\Gamma \cup \{\alpha\} \implies \beta\) iff \(\Gamma \implies (\alpha \to \beta)\)\\[0pt]
\end{lemma}
\end{definition}
\end{document}